
\documentclass[12pt]{article}
\usepackage{amsmath, amssymb, amsthm}
\usepackage{geometry}
\usepackage{hyperref}
\geometry{margin=1in}
\title{Toward a Resolution of the Hodge Conjecture via Algebraic Convergence}
\author{Thomas Birnie (in collaboration with ChatGPT)}
\date{}

\begin{document}
\maketitle

\section*{Abstract}
We propose a strategy for resolving the Hodge Conjecture by constructing rational Hodge classes as limits of sequences of algebraic cycle classes. This framework incorporates tools from sheaf cohomology, derived categories, and algebraic topology, and reframes transcendental obstructions as convergent anomalies in Mayer–Vietoris constructions. We outline a pathway from approximable cycles to canonical representatives through a derived-categorical and motivic lifting process.

\section{Problem Statement}
Let \( X \) be a smooth projective complex variety. The classical Hodge Conjecture asserts:
\begin{quote}
Every class \( \gamma \in H^{2p}(X, \mathbb{Q}) \cap H^{p,p}(X) \) is the image of an algebraic cycle class:
\[
\exists Z \in Z^p(X): \quad \operatorname{cl}(Z) = \gamma.
\]
\end{quote}

In known cases, obstructions to constructing such \( Z \) often appear in higher cohomology groups, particularly as connecting morphisms in long exact sequences:
\[
\delta(\alpha_U - \alpha_V) = \tau \in H^{2p+1}(X).
\]

\section{Core Proposal}
We replace the direct construction with a convergence framework:
\begin{enumerate}
    \item \textbf{Construct a sequence of algebraic cycles} \( Z_n \in Z^p(X) \) such that
    \[
    \lim_{n \to \infty} \operatorname{cl}(Z_n) = \gamma_\infty \in H^{2p}(X, \mathbb{Q}) \cap H^{p,p}(X).
    \]
    \item \textbf{Interpret gluing mismatches} \( \tau_n = \delta(\alpha_U^{(n)} - \alpha_V^{(n)}) \in H^{2p+1}(X) \) as a converging error term:
    \[
    \lim_{n \to \infty} \tau_n = \tau.
    \]
    \item \textbf{Lift the obstruction} \( \tau \) to a trivial class in a derived category or via a motivic extension, such that it becomes exact or cancelable.
    \item \textbf{Compactify the sequence} \( \{Z_n\} \) into a single algebraic cycle \( Z \) through a convergence-preserving functor from a moduli space, stack, or formal scheme.
\end{enumerate}

\section{Strategic Hypothesis}
We hypothesize:

\begin{quote}
\textbf{Algebraic Approximability of Hodge Classes:} Every rational Hodge class is the limit of a sequence of algebraic cycle classes.
\end{quote}

\section{Technical Goals}
\begin{itemize}
    \item Define and prove convergence of \( \{Z_n\} \) under a Hodge-norm topology.
    \item Construct a categorical compactification functor \( \widehat{\operatorname{cl}}: \widehat{Z}^p(X) \to H^{2p}(X, \mathbb{Q}) \).
    \item Show that the limiting class \( \gamma_\infty \) corresponds to an actual cycle via patching, degeneration, or formal limit stabilization.
    \item Show that all \( \tau \in H^{2p+1}(X) \) arising from such sequences are killed in a derived or motivic category.
\end{itemize}

\section{Conclusion}
If successful, this program will convert a transcendental obstruction into a geometric convergence process, yielding a genuine algebraic cycle representative for every rational Hodge class and completing the classical conjecture.

\end{document}
